%!TEX TS-program = pdflatex

\documentclass[aspectratio=169,10pt]{beamer}

\usepackage[T2A]{fontenc}
\usepackage[utf8]{inputenc}
\usepackage[english,russian]{babel}
\usepackage{cmap}
\usepackage{amsmath,amssymb,mathtools}
\usepackage{booktabs}
\usepackage{array}
\usepackage{etoolbox}
\usepackage{PTSans}
\usepackage{graphicx}
\graphicspath{{figures/}}

\renewcommand{\familydefault}{\sfdefault}
\renewcommand{\rmdefault}{\sfdefault}

\newbool{russian}
\booltrue{russian}
\usepackage{HSE-theme/beamerthemeHSE}

\hypersetup{bookmarks=false}
\setbeamertemplate{blocks}[rounded][shadow=false]
\setbeamerfont{frametitle}{size=\Large}
\setbeamerfont{framesubtitle}{size=\small}
\setbeamerfont{block title}{size=\normalsize}
\setbeamerfont{block body}{size=\small}

\newcommand{\R}{\mathbb{R}}
\newcommand{\norm}[1]{\left\lVert #1\right\rVert}
\newcommand{\grad}{\nabla}
\newcommand{\argmin}{\operatorname*{arg\,min}}
\newcommand{\lammin}{\lambda_{\min}}
\newcommand{\lammax}{\lambda_{\max}}
\newcommand{\kappaA}{\kappa(A)}

\title[Методы оптимизации в ML]{Методы оптимизации\\в машинном обучении}
\subtitle{Лекция 4. Обусловленность и градиентный спуск: почему алгоритм сходится?}
\author[Н. Лукьяненко]{Никита Лукьяненко}
\institute{Совместная программа СМОЛ ГУ и ФКН НИУ ВШЭ}
\date{2026/27 учебный год}

\begin{document}

% 1
\frame[plain]{\titlepage}

% 2
\begin{frame}{С чего продолжаем}
\small
На прошлой лекции для квадратичной функции
\[
 f(x)=\frac12x^\top Ax-b^\top x+c,
 \qquad A=A^\top\succ0,
\]
мы ввели число обусловленности
\[
 \boxed{\kappaA=\frac{\lammax(A)}{\lammin(A)}\ge1.}
\]

\begin{alertblock}{Сегодня}
Сначала разберёмся, что именно измеряет $\kappa(A)$ и почему эта величина важна. Затем увидим, как тот же спектр определяет допустимый шаг и скорость градиентного спуска.
\end{alertblock}
\end{frame}

% 3
\begin{frame}{Геометрический смысл числа обусловленности}
\centering
\includegraphics[width=0.94\linewidth]{conditioning_contours.pdf}

\vspace{0.2em}
\small
$\kappa$ измеряет не общую ``крутизну'' функции, а \textbf{неравномерность кривизны по разным направлениям}.
\end{frame}

% 4
\begin{frame}{Экстремальные кривизны}
\small
Для симметричной положительно определённой матрицы
\[
 \boxed{
 \min_{\norm{v}=1}v^\top Av=\lammin(A),
 \qquad
 \max_{\norm{v}=1}v^\top Av=\lammax(A).
 }
\]

Поэтому
\[
 \boxed{
 \kappaA=
 \frac{\displaystyle\max_{\norm{v}=1}v^\top Av}
      {\displaystyle\min_{\norm{v}=1}v^\top Av}.
 }
\]

\begin{block}{Интерпретация}
Число обусловленности показывает, во сколько раз самая большая направленная кривизна превосходит самую маленькую.
\end{block}
\end{frame}

% 5
\begin{frame}{Почему вытянутость равна $\sqrt{\kappa}$, а не $\kappa$?}
\footnotesize
В собственных координатах линия уровня квадратичной функции имеет вид
\[
 \frac12\sum_i\lambda_i z_i^2=c.
\]
Вдоль собственного направления $v_i$ остальные координаты равны нулю, поэтому
\[
 \frac12\lambda_i z_i^2=c
 \quad\Longrightarrow\quad
 |z_i|=a_i=\sqrt{\frac{2c}{\lambda_i}}.
\]
Значит,
\[
 \boxed{
 \frac{a_{\max}}{a_{\min}}
 =\sqrt{\frac{\lammax}{\lammin}}
 =\sqrt{\kappaA}.
 }
\]

\begin{alertblock}{Пример}
Если $\kappa=100$, линия уровня вытянута по главным осям в $10$ раз.
\end{alertblock}
\end{frame}

% 6
\begin{frame}{Масштаб функции и обусловленность --- разные вещи}
\small
Если умножить матрицу на $c>0$, все собственные значения умножатся на $c$:
\[
 \lambda_i(cA)=c\lambda_i(A).
\]
Следовательно,
\[
 \boxed{
 \kappa(cA)
 =\frac{c\lammax(A)}{c\lammin(A)}
 =\kappa(A).
 }
\]

\begin{block}{Самопроверка}
Что хуже обусловлено?
\[
 A_1=0.01I,
 \qquad
 A_2=\operatorname{diag}(1,100).
\]
\end{block}

Ответ: $\kappa(A_1)=1$, а $\kappa(A_2)=100$.
\end{frame}

% 7
\begin{frame}{Почему это вообще называется ``обусловленностью''?}
\small
Рассмотрим линейную систему
\[
 Ax=b,
 \qquad A\succ0.
\]
Её решение
\[
 x=A^{-1}b
\]
зависит от входных данных $b$.

\begin{block}{Вопрос}
Если слегка изменить $b$, насколько сильно может измениться решение $x$?
\end{block}

Именно эту \textbf{чувствительность задачи} контролирует число обусловленности.
\end{frame}

% 8
\begin{frame}{Одинаковое возмущение --- разный сдвиг решения}
\small
Рассмотрим одно и то же относительное возмущение правой части:
\[
 b=\begin{pmatrix}1\\0\end{pmatrix},\qquad
 \Delta b=\begin{pmatrix}0\\0.1\end{pmatrix},\qquad
 \frac{\norm{\Delta b}}{\norm b}=10\%.
\]

\vspace{0.15em}
\centering
\includegraphics[width=0.82\linewidth]{sensitivity_solution.pdf}

\vspace{0.15em}
\footnotesize
Одинаковое изменение входных данных может почти не изменить решение хорошо обусловленной системы и сильно изменить решение плохо обусловленной.
\end{frame}

% 9
\begin{frame}{Оценка чувствительности решения}
\small
\begin{block}{Утверждение}
Пусть $A=A^\top\succ0$, $Ax=b$, а при фиксированной $A$ правая часть получила возмущение $\Delta b$. Тогда при $b\neq0$
\[
 \boxed{
 \frac{\norm{\Delta x}}{\norm{x}}
 \le
 \kappaA\,
 \frac{\norm{\Delta b}}{\norm{b}}.
 }
\]
\end{block}

\begin{alertblock}{Смысл}
$\kappa(A)$ --- максимально возможный коэффициент усиления относительной ошибки правой части в этой оценке.
\end{alertblock}
\end{frame}

% 10
\begin{frame}{Доказательство оценки чувствительности: первая половина}
\footnotesize
После возмущения имеем
\[
 A(x+\Delta x)=b+\Delta b.
\]
Вычитая $Ax=b$, получаем
\[
 A\Delta x=\Delta b,
 \qquad
 \Delta x=A^{-1}\Delta b.
\]
Собственные значения $A^{-1}$ равны $1/\lambda_i(A)$, поэтому
\[
 \norm{A^{-1}}_2=\frac1{\lammin(A)}.
\]
Следовательно,
\[
 \boxed{
 \norm{\Delta x}
 \le
 \frac1{\lammin(A)}\norm{\Delta b}.
 }
\]
\end{frame}

% 11
\begin{frame}{Доказательство оценки чувствительности: завершаем}
\footnotesize
Так как $b=Ax$,
\[
 \norm{b}=\norm{Ax}\le\lammax(A)\norm{x}.
\]
Отсюда
\[
 \frac1{\norm{x}}
 \le
 \frac{\lammax(A)}{\norm{b}}.
\]
Комбинируем с предыдущей оценкой:
\[
 \frac{\norm{\Delta x}}{\norm{x}}
 \le
 \frac{\lammax(A)}{\lammin(A)}
 \frac{\norm{\Delta b}}{\norm{b}}.
\]
То есть
\[
 \boxed{
 \frac{\norm{\Delta x}}{\norm{x}}
 \le
 \kappaA\frac{\norm{\Delta b}}{\norm{b}}.
 }
\]
\end{frame}

% 12
\begin{frame}{Обусловленность задачи и скорость алгоритма --- не одно и то же}
\small
До сих пор мы говорили о \textbf{чувствительности самой задачи}: как решение реагирует на изменение входных данных.

\vspace{0.6em}
Теперь переключимся на \textbf{алгоритм}: сколько шагов потребуется, чтобы добраться до решения?

\begin{alertblock}{Интрига}
Для квадратичных задач оказывается, что те же числа $\lammin$, $\lammax$ и $\kappa$ управляют и поведением градиентного спуска.
\end{alertblock}
\end{frame}

% 13
\begin{frame}{Напоминание: градиентный спуск}
\small
В курсе машинного обучения вы уже встречали алгоритм
\[
 \boxed{
 x_{k+1}=x_k-\alpha\grad f(x_k),
 }
\]
где $\alpha>0$ --- шаг.

\begin{block}{Сегодня нас интересуют не вычислительные детали}
Мы хотим математически ответить на два вопроса:
\begin{enumerate}
 \item почему этот шаг вообще должен уменьшать функцию;
 \item при каких $\alpha$ последовательность действительно сходится.
\end{enumerate}
\end{block}
\end{frame}

% 14
\begin{frame}{Откуда возникает направление $-\grad f(x)$?}
\small
Из второй лекции:
\[
 f(x+h)\approx f(x)+\grad f(x)^\top h.
\]
Если выбирать направление единичной длины, максимально быстро уменьшает линейную модель направление
\[
 -\frac{\grad f(x)}{\norm{\grad f(x)}}.
\]

Но одной линейной модели недостаточно для выбора \textbf{длины} шага.
\end{frame}

% 15
\begin{frame}{Линейная модель сама не умеет выбирать длину шага}
\small
Локально мы заменяем исходную функцию линейной моделью
\[
 m_{\rm lin}(h)=\grad f(x)^\top h.
\]
Если $\grad f(x)\neq0$ и $h=-t\grad f(x)$, $t>0$, то
\[
 m_{\rm lin}(-t\grad f(x))=-t\norm{\grad f(x)}^2\xrightarrow[t\to\infty]{}-\infty.
\]

\begin{alertblock}{Важно}
Неограниченно уменьшается именно \textbf{линейная модель}, а не исходная функция $f$. Линейное приближение надёжно только для малых $\norm h$, поэтому само по себе оно не может выбрать длину шага.
\end{alertblock}

\centering
\includegraphics[width=0.67\linewidth]{regularized_local_model.pdf}
\end{frame}

% 16
\begin{frame}{Как найти минимум вспомогательной модели?}
\small
Рассмотрим функцию шага
\[
 m_x(h)=\grad f(x)^\top h+\frac1{2\alpha}\norm h^2,
 \qquad \alpha>0.
\]

\begin{columns}[T,onlytextwidth]
\column{0.48\textwidth}
\begin{block}{1. Необходимое условие}
Если дифференцируемая функция имеет внутренний локальный минимум в $h^*$, то
\[
 \grad_h m_x(h^*)=0.
\]
Здесь
\[
 \grad_hm_x(h)=\grad f(x)+\frac1\alpha h.
\]
\end{block}

\column{0.48\textwidth}
\begin{block}{2. Почему здесь этого достаточно?}
Гессиан по переменной $h$ равен
\[
 \nabla_h^2m_x(h)=\frac1\alpha I\succ0.
\]
Поэтому $m_x$ --- строго выпуклая квадратичная функция: её стационарная точка единственна и является глобальным минимумом.
\end{block}
\end{columns}

\vspace{0.35em}
\centering
Следовательно, чтобы найти $h^*$, достаточно решить уравнение $\grad_hm_x(h)=0$.
\end{frame}

% 17
\begin{frame}{Градиентный шаг как решение локальной задачи}
\small
Из предыдущего уравнения
\[
 \boxed{h^*=-\alpha\grad f(x).}
\]
Поэтому
\[
 x_{k+1}=x_k+h^*
 =\boxed{x_k-\alpha\grad f(x_k)}.
\]

\begin{alertblock}{Важно}
Мы вывели форму шага, но пока \textbf{не доказали}, что исходная функция $f$ при этом уменьшается или что последовательность сходится.
\end{alertblock}
\end{frame}

% 18
\begin{frame}{Главный вопрос: как выбрать $\alpha$?}
\small
Слишком маленький шаг может дать очень медленное движение.

Слишком большой --- перескакивание через минимум и даже расходимость.

\begin{center}
\Large
\[
 \boxed{\text{Какие значения }\alpha\text{ гарантируют сходимость?}}
\]
\end{center}

Начнём с квадратичной функции, где ответ можно получить полностью.
\end{frame}

% 19
\begin{frame}{Квадратичная задача: всё уже известно}
\small
Пусть
\[
 f(x)=\frac12x^\top Ax-b^\top x+c,
 \qquad A=A^\top\succ0.
\]
Из третьей лекции:
\[
 \grad f(x)=Ax-b,
 \qquad
 x^*=A^{-1}b.
\]
Градиентный шаг принимает вид
\[
 \boxed{
 x_{k+1}=x_k-\alpha(Ax_k-b).
 }
\]
\end{frame}

% 20
\begin{frame}{Вместо самих точек будем изучать ошибку}
\small
Определим
\[
 \boxed{e_k=x_k-x^*.}
\]
Тогда сходимость
\[
 x_k\to x^*
\]
эквивалентна
\[
 e_k\to0.
\]

\begin{block}{Почему это удобно?}
Уравнение $Ax^*=b$ позволит убрать из рекурсии линейный член $b$.
\end{block}
\end{frame}

% 21
\begin{frame}{Получаем линейную рекурсию для ошибки}
\footnotesize
\[
\begin{aligned}
 e_{k+1}
 &=x_{k+1}-x^*\\
 &=x_k-\alpha(Ax_k-b)-x^*.
\end{aligned}
\]
Так как $b=Ax^*$,
\[
 e_{k+1}
 =x_k-x^*-\alpha A(x_k-x^*).
\]
Следовательно,
\[
 \boxed{
 e_{k+1}=(I-\alpha A)e_k.
 }
\]

Теперь задача сходимости свелась к анализу степеней одной матрицы.
\end{frame}

% 22
\begin{frame}{Переходим в собственный базис}
\small
Поскольку $A$ симметрична, существует ортонормированный базис собственных векторов $v_i$:
\[
 Av_i=\lambda_i v_i.
\]
Разложим ошибку:
\[
 e_k=\sum_i z_{k,i}v_i.
\]
Тогда
\[
 (I-\alpha A)v_i
 =(1-\alpha\lambda_i)v_i.
\]

\begin{alertblock}{Главная идея}
В собственных координатах многомерный градиентный спуск распадается на независимые одномерные рекурсии.
\end{alertblock}
\end{frame}

% 23
\begin{frame}{Рекурсия для каждой компоненты}
\small
Подставляя разложение в
\[
 e_{k+1}=(I-\alpha A)e_k,
\]
получаем
\[
 \boxed{
 z_{k+1,i}
 =(1-\alpha\lambda_i)z_{k,i}.
 }
\]

Введём обозначение
\[
 \boxed{q_i:=1-\alpha\lambda_i.}
\]
Тогда
\[
 z_{k+1,i}=q_i z_{k,i},
 \qquad
 z_{k,i}=q_i^k z_{0,i}.
\]

\begin{alertblock}{Смысл $q_i$}
$q_i$ --- множитель, на который меняется $i$-я компонента ошибки за одну итерацию.
\end{alertblock}
\end{frame}

% 24
\begin{frame}{Одна координата: три возможных режима}
\centering
\includegraphics[width=0.95\linewidth]{scalar_recurrence_regimes.pdf}

\vspace{0.15em}
\small
Знак $q_i$ определяет наличие колебаний, а модуль $|q_i|$ --- затухает ошибка или растёт.
\end{frame}

% 25
\begin{frame}{Условие сходимости в одном собственном направлении}
\small
Для рекурсии
\[
 z_{k+1}=qz_k
\]
имеем $z_k\to0$ при любом $z_0$ тогда и только тогда, когда
\[
 |q|<1.
\]
Здесь
\[
 q=1-\alpha\lambda.
\]
Поэтому требуется
\[
 |1-\alpha\lambda|<1.
\]
\end{frame}

% 26
\begin{frame}{Что означает $|1-\alpha\lambda|<1$?}
\centering
\includegraphics[width=0.76\linewidth]{stability_interval.pdf}

\vspace{0.1em}
\small
На рисунке взято $\lambda=3$. Поэтому устойчивый интервал имеет вид
\[
0<\alpha<\frac{2}{3}.
\]
В точке $\alpha=1/3$ получаем $q=0$: эта компонента ошибки исчезает за один шаг.
\end{frame}

% 27
\begin{frame}{Решаем неравенство}
\small
Так как $\lambda>0$,
\[
 |1-\alpha\lambda|<1
\]
равносильно
\[
 -1<1-\alpha\lambda<1.
\]
Отсюда
\[
 0<\alpha\lambda<2,
\]
то есть
\[
 \boxed{
 0<\alpha<\frac2\lambda.
 }
\]

Для многомерной задачи условие должно выполняться для \textbf{каждого} собственного значения.
\end{frame}

% 28
\begin{frame}{Теорема о сходимости на квадратичной функции}
\small
\begin{block}{Теорема}
Пусть $A=A^\top\succ0$. Тогда метод
\[
 x_{k+1}=x_k-\alpha(Ax_k-b)
\]
сходится к $x^*=A^{-1}b$ при любом $x_0$ тогда и только тогда, когда
\[
 \boxed{
 0<\alpha<\frac{2}{\lammax(A)}.
 }
\]
\end{block}

Это первый полный ответ на вопрос о выборе шага.
\end{frame}

% 29
\begin{frame}{Доказательство теоремы}
\footnotesize
Мы получили
\[
 z_{k,i}=(1-\alpha\lambda_i)^kz_{0,i}.
\]
Поэтому $e_k\to0$ при любом начальном $e_0$ тогда и только тогда, когда
\[
 |1-\alpha\lambda_i|<1
 \qquad\forall i.
\]
Для каждого $i$ это равносильно
\[
 0<\alpha<\frac{2}{\lambda_i}.
\]
Чтобы условие выполнялось одновременно для всех $i$, достаточно и необходимо взять самое строгое ограничение:
\[
 0<\alpha<\min_i\frac2{\lambda_i}
 =\frac2{\lammax(A)}.
\]
Следовательно, $e_k\to0$ и $x_k\to x^*$. \hfill$\square$
\end{frame}

% 30
\begin{frame}{Не только сходимость: насколько быстро уменьшается ошибка?}
\small
Из собственных координат
\[
 |z_{k,i}|
 =|1-\alpha\lambda_i|^k|z_{0,i}|.
\]
Определим
\[
 \boxed{
 q(\alpha)=\max_i|1-\alpha\lambda_i|.
 }
\]
Тогда
\[
 \boxed{
 \norm{e_k}\le q(\alpha)^k\norm{e_0}.
 }
\]

\begin{block}{Интерпретация}
Чем меньше $q$, тем быстрее геометрически убывает ошибка.
\end{block}
\end{frame}

% 31
\begin{frame}{Почему при $\alpha=1/L$ худшее направление соответствует $\mu$?}
\small
Обозначим $L=\lammax(A)$ и $\mu=\lammin(A)$. Для каждой собственной компоненты ошибки после одного шага остаётся доля
\[
 q_i=\left|1-\frac{\lambda_i}{L}\right|
 =1-\frac{\lambda_i}{L},
 \qquad \lambda_i\in[\mu,L].
\]

\begin{columns}[T,onlytextwidth]
\column{0.50\textwidth}
\begin{block}{Что значит ``худшее''?}
Чем \textbf{больше} $q_i$, тем большая доля ошибки остаётся после шага и тем медленнее затухает эта компонента.

Поскольку $q_i$ убывает с ростом $\lambda_i$, максимум достигается при $\lambda_i=\mu$:
\[
 q_{\max}=1-\frac\mu L
 =1-\frac1\kappa.
\]
\end{block}
\column{0.46\textwidth}
\centering
\includegraphics[width=\linewidth]{worst_direction_factor.pdf}
\end{columns}

\vspace{0.1em}
\footnotesize
То есть $\mu/L$ --- доля, которая гарантированно ``снимается'' в самом медленном направлении; $1-\mu/L$ --- доля ошибки, которая остаётся.
\end{frame}

% 32
\begin{frame}{Почему большое $\kappa$ означает медленную сходимость?}
\centering
\includegraphics[width=0.78\linewidth]{contraction_vs_kappa.pdf}
\end{frame}

% 33
\begin{frame}{Геометрия траектории при разной обусловленности}
\centering
\includegraphics[width=0.95\linewidth]{gd_condition_paths.pdf}
\end{frame}

% 34
\begin{frame}{Почему обусловленность замедляет градиентный спуск?}
\small
\vspace{0.55em}
При шаге $\alpha=1/L$ мы получили оценку
\[
 \boxed{
 \norm{e_k}
 \le
 \left(1-\frac1\kappa\right)^k\norm{e_0}.
 }
\]

\centering
\includegraphics[width=0.58\linewidth]{error_decay_conditioning.pdf}

\vspace{0.05em}
\begin{alertblock}{Вывод}
При большом $\kappa$ множитель $1-1/\kappa$ близок к единице, поэтому гарантированное сокращение ошибки может быть очень медленным.
\end{alertblock}
\end{frame}

% 35
\begin{frame}{Но реальные функции не обязаны быть квадратичными}
\small
Для квадратичной функции
\[
 \grad^2 f(x)=A
\]
постоянен.

Для общей гладкой функции кривизна меняется от точки к точке.

\begin{alertblock}{Новая задача}
Нужна характеристика, которая позволит контролировать, насколько быстро может меняться градиент.
\end{alertblock}
\end{frame}

% 36
\begin{frame}{Почему величины градиента в одной точке недостаточно}
\centering
\includegraphics[width=0.88\linewidth]{gradient_change_nonquadratic.pdf}
\end{frame}

% 37
\begin{frame}{Липшицевость: общее определение}
\small
\begin{block}{Определение}
Отображение $g:\R^d\to\R^m$ называется \textbf{липшицевым с константой $L\ge0$}, если для любых $x,y$
\[
 \boxed{
 \norm{g(x)-g(y)}\le L\norm{x-y}.
 }
\]
\end{block}

\begin{alertblock}{Что означает это неравенство?}
Изменение значения $g$ не может быть больше, чем $L$ раз изменение аргумента. Константа $L$ задаёт глобальное ограничение на скорость изменения отображения.
\end{alertblock}

Если условие выполняется с некоторым $L$, оно автоматически выполняется и с любой большей константой.
\end{frame}

% 38
\begin{frame}{$L$-гладкость}
\small
\begin{block}{Определение}
Дифференцируемая функция $f:\R^d\to\R$ называется $L$-гладкой, если её градиент $\grad f$ является липшицевым с константой $L$:
\[
 \boxed{
 \norm{\grad f(x)-\grad f(y)}
 \le L\norm{x-y}.
 }
\]
\end{block}

\begin{alertblock}{Смысл}
$L$ ограничивает скорость изменения \textbf{градиента}, а не самой функции. Большое $L$ означает, что направление и величина градиента могут быстро меняться даже на небольших расстояниях.
\end{alertblock}
\end{frame}

% 38
\begin{frame}{Как понимать параметр $L$ в одном измерении?}
\small
Пусть $g(t)=f'(t)$. Липшицевость градиента означает, что для любых $t_1\neq t_2$
\[
 \boxed{\frac{|g(t_2)-g(t_1)|}{|t_2-t_1|}\le L.}
\]
То есть $L$ ограничивает \textbf{наклон любой секущей} графика производной $g=f'$.

\vspace{0.15em}
\centering
\includegraphics[width=0.78\linewidth]{smoothness_gradient_change.pdf}

\vspace{0.1em}
\footnotesize
Если $f''$ непрерывна, локально это соответствует условию $|f''(t)|\le L$: гладкость ограничивает величину кривизны.
\end{frame}

% 39
\begin{frame}{Квадратичная функция как проверка определения}
\small
Пусть
\[
 f(x)=\frac12x^\top Ax-b^\top x,
 \qquad A=A^\top\succeq0.
\]
Тогда
\[
 \grad f(x)-\grad f(y)=A(x-y).
\]
Следовательно,
\[
 \norm{\grad f(x)-\grad f(y)}
 \le \norm{A}_2\norm{x-y}.
\]
Для симметричной положительно полуопределённой $A$
\[
 \norm{A}_2=\lammax(A).
\]
Значит, квадратичная функция $L$-гладкая при
\[
 \boxed{L=\lammax(A).}
\]
\end{frame}

% 40
\begin{frame}{Гладкость даёт квадратичную верхнюю модель}
\small
\begin{block}{Лемма об убывании}
Если $f$ имеет $L$-липшицев градиент, то для любых $x,h$
\[
 \boxed{
 f(x+h)
 \le
 f(x)+\grad f(x)^\top h
 +\frac L2\norm{h}^2.
 }
\]
\end{block}

Для квадратичной функции такая формула была точной. Для общей гладкой функции она становится \textbf{глобальной верхней оценкой}.
\end{frame}

% 41
\begin{frame}{Геометрия леммы об убывании}
\centering
\includegraphics[width=0.83\linewidth]{smooth_upper_bound.pdf}
\end{frame}

% 42
\begin{frame}{Доказательство леммы: сводим всё к одной переменной}
\footnotesize
Рассмотрим путь
\[
 \varphi(t)=f(x+th),
 \qquad t\in[0,1].
\]
По основной теореме анализа
\[
 f(x+h)-f(x)
 =\int_0^1\varphi'(t)\,dt.
\]
По правилу цепочки
\[
 \varphi'(t)=\grad f(x+th)^\top h.
\]
Следовательно,
\[
 f(x+h)-f(x)
 =\int_0^1\grad f(x+th)^\top h\,dt.
\]
\end{frame}

% 43
\begin{frame}{Доказательство леммы: используем гладкость}
\footnotesize
Вычтем линейную часть:
\[
\begin{aligned}
&f(x+h)-f(x)-\grad f(x)^\top h\\
&=\int_0^1
[\grad f(x+th)-\grad f(x)]^\top h\,dt.
\end{aligned}
\]
По неравенству Коши--Буняковского и $L$-гладкости
\[
 [\grad f(x+th)-\grad f(x)]^\top h
 \le Lt\norm{h}^2.
\]
Интегрируя,
\[
 \int_0^1Lt\norm{h}^2dt
 =\frac L2\norm{h}^2.
\]
Получаем требуемую оценку. \hfill$\square$
\end{frame}

% 44
\begin{frame}{Подставляем градиентный шаг}
\small
Положим
\[
 h=-\alpha\grad f(x).
\]
Тогда из леммы
\[
\begin{aligned}
 f(x-\alpha\grad f(x))
 &\le f(x)-\alpha\norm{\grad f(x)}^2\\
 &\quad+\frac L2\alpha^2\norm{\grad f(x)}^2.
\end{aligned}
\]
То есть
\[
 \boxed{
 f(x_{k+1})
 \le
 f(x_k)
 -\alpha\left(1-\frac{L\alpha}{2}\right)
 \norm{\grad f(x_k)}^2.
 }
\]
\end{frame}

% 45
\begin{frame}{Теорема об убывании градиентного спуска}
\small
\begin{block}{Следствие}
Если $f$ $L$-гладкая и
\[
 0<\alpha<\frac2L,
\]
то при $\grad f(x_k)\neq0$
\[
 f(x_{k+1})<f(x_k).
\]
\end{block}

Особенно удобно взять
\[
 \alpha=\frac1L.
\]
Тогда
\[
 \boxed{
 f(x_{k+1})
 \le
 f(x_k)-\frac1{2L}\norm{\grad f(x_k)}^2.
 }
\]
\end{frame}

% 46
\begin{frame}{Как выглядит гарантированное убывание}
\centering
\includegraphics[width=0.80\linewidth]{descent_step_visual.pdf}
\end{frame}

% 47
\begin{frame}{Что мы доказали --- и чего ещё не доказали}
\small
Из $L$-гладкости мы получили:
\[
 f(x_{k+1})\le f(x_k)
\]
при подходящем шаге.

Но отсюда \textbf{не следует}, что
\[
 x_k\to x^*_{\mathrm{global}}.
\]

\begin{block}{Почему?}
Невыпуклая функция может иметь локальные минимумы и седловые точки. Убывание функции --- это ещё не глобальная оптимальность.
\end{block}

Для более сильных выводов нам понадобятся дополнительные свойства функции.
\end{frame}

% 48
\begin{frame}{Как влияет размер шага на траекторию?}